% ==============================================================================
% CAPÍTULO 2: FUNDAMENTAÇÃO TEÓRICA
%
% EQUAÇÕES NESTE CAPÍTULO: apenas as que DEFINEM um conceito, modelo ou teoria
% da literatura ("o que é" / "como a literatura o formaliza"), sempre com a
% citação da fonte original. Métricas, parâmetros e procedimentos de cálculo
% escolhidos pelo pesquisador ficam na Metodologia, que remete às equações
% daqui com \eqref, sem repeti-las. Ver o exemplo na Seção 2.1.
% Os blocos \begin{orientacao}...\end{orientacao} aparecem no PDF apenas no
% modo revisão e somem no modo final; ainda assim, apague-os após a leitura.
% ==============================================================================
\chapter{Fundamentação teórica}
\label{ch:fundamentacao_teorica}
\begin{resumocapitulo}
  Aqui deve ser incluído um resumo do capítulo, destacando os principais conceitos, teorias e abordagens que fundamentam a pesquisa. O resumo deve apresentar de forma concisa os tópicos abordados, a relevância teórica e metodológica das discussões e como elas se conectam com os objetivos do trabalho. Deve ser escrito em parágrafo único, sem citações ou referências bibliográficas e de forma concisa e não deve ultrapassar 200 palavras.
\end{resumocapitulo}

A fundamentação teórica tem como finalidade articular e sintetizar as principais correntes conceituais, modelos analíticos e construtos que sustentam a investigação. O seu papel estrutural é fornecer o alicerce epistemológico necessário para a compreensão do objeto de estudo, estabelecendo um diálogo crítico com a literatura consolidada e o estado da arte. Ela define a lente teórica pela qual os fenômenos serão observados e os dados empíricos, interpretados. As fontes discutidas neste capítulo foram recuperadas segundo o protocolo descrito na Seção~\ref{sec:procedimentos_literatura}.

\section{Construto teórico 1}
\label{sec:construto_1}
Cada seção conceitual deve apresentar um construto central da pesquisa, sua definição operacional, as principais vertentes teóricas que o discutem e os pontos de convergência ou tensão entre autores. Não se trata de acumular definições, mas de posicionar o estudo diante delas, explicitando qual entendimento será adotado e por quê.

% ------------------------------------------------------------------------------
% EXEMPLO: equação conceitual na fundamentação teórica (adaptar ou apagar)
% Função da equação aqui: DEFINIR o conceito, como formalizado na literatura.
% Por isso: (1) vem com a citação da fonte original; (2) usa a notação do
% autor de origem; (3) NÃO traz parâmetros, bibliotecas ou valores escolhidos
% neste estudo. A Metodologia (Seção 3.6) retoma esta equação por \eqref.
% ------------------------------------------------------------------------------
\begin{orientacao}[Equações na fundamentação teórica]
  Neste capítulo entram apenas as expressões que definem formalmente um conceito, modelo ou teoria da literatura, isto é, as que respondem à pergunta \emph{o que é} e \emph{como a literatura o formaliza}. Cada uma deve vir acompanhada da citação da fonte original. Métricas de avaliação, parâmetros escolhidos e procedimentos de cálculo pertencem à Metodologia (Capítulo~\ref{ch:metodologia}), que deve remeter à equação apresentada aqui, sem repeti-la. O exemplo a seguir ilustra o caso: a definição dos valores de Shapley aparece apenas neste capítulo, e a Seção~\ref{sec:analise} descreve como eles foram calculados no estudo.
\end{orientacao}

Entre as abordagens de explicabilidade de modelos, destaca-se a atribuição de importância fundamentada na teoria dos jogos cooperativos. \citeonline{Shapley1953} propôs um critério de repartição do ganho de uma coalizão entre seus participantes, segundo o qual a contribuição do jogador $j$ corresponde à média ponderada de suas contribuições marginais a todas as coalizões possíveis, conforme a Equação~\eqref{eq:shapley}:
\begin{equation}
  \phi_j(v) = \sum_{S \subseteq N \setminus \{j\}} \frac{|S|!\,\left(|N| - |S| - 1\right)!}{|N|!}\,\Big[ v\left(S \cup \{j\}\right) - v(S) \Big],
  \label{eq:shapley}
\end{equation}
em que $N$ é o conjunto de jogadores, $S$ é uma coalizão que não contém $j$ e $v$ é a função característica, que associa a cada coalizão o ganho por ela obtido. \citeonline{Lundberg2017} transpuseram essa formulação para a explicação de modelos de aprendizado de máquina, tratando as variáveis preditoras como jogadores e a predição do modelo como o ganho a ser repartido, o que deu origem ao método SHAP (\textit{SHapley Additive exPlanations}).

\section{Construto teórico 2}
\label{sec:construto_2}
A segunda seção conceitual segue a mesma lógica. Recomenda-se que a sequência dos construtos acompanhe a ordem dos objetivos específicos, o que facilita ao leitor perceber como cada base teórica será mobilizada nas etapas analíticas.

A seguir são apresentados exemplos de utilização de hierarquias de seções

% ------------------------------------------------------------------------------
% EXEMPLO: hierarquia de seções (adaptar ou apagar)
% Níveis disponíveis (numeração progressiva, NBR 6024):
%   \chapter        1        primária
%   \section        1.1      secundária
%   \subsection     1.1.1    terciária
%   \subsubsection  1.1.1.1  quaternária
% Cada título recebe um \label LOGO APÓS o comando, com prefixo que indica
% o nível e nome curto, sem acentos, espaços ou maiúsculas:
%   sec:nome   subsec:nome   subsubsec:nome
% Remissão no texto: Seção~\ref{subsec:nome} (o ~ impede a quebra de linha
% entre a palavra e o número).
% ------------------------------------------------------------------------------
\subsection{Esta é uma subseção}
\label{subsec:exemplo_subsecao}

\begin{orientacao}[Seções, subseções e rótulos]
  A subseção (\texttt{\textbackslash subsection}) divide uma seção em partes com unidade temática própria, e a subsubseção (\texttt{\textbackslash subsubsection}) detalha um aspecto específico de uma subseção. Escreva os títulos apenas com a inicial maiúscula, pois a formatação é aplicada pelo modelo. Abra uma subdivisão somente se houver pelo menos duas no mesmo nível: uma subseção 2.3.1 isolada, sem 2.3.2, indica que o conteúdo pertence à própria seção 2.3. Evite também dois títulos seguidos sem texto entre eles; cada nível deve começar com ao menos um parágrafo que anuncie o que vem a seguir. Para remeter a uma divisão, use \texttt{Seção\textasciitilde\textbackslash ref\{rótulo\}}: o número é atualizado automaticamente se a ordem do texto mudar. Prefira não passar do quarto nível de numeração, pois divisões muito profundas dificultam a leitura.
\end{orientacao}

Texto introdutório da subseção, que apresenta o conteúdo das subsubseções seguintes. A Seção~\ref{subsubsec:exemplo_subsubsecao} detalha o primeiro aspecto.

\subsubsection{Esta é uma subsubseção}
\label{subsubsec:exemplo_subsubsecao}

Texto da subsubseção. Para remeter à subseção que a contém, escreva Seção~\ref{subsec:exemplo_subsecao}.

\section{Trabalhos relacionados}
\label{sec:trabalhos_relacionados}
A seção de trabalhos relacionados tem como finalidade examinar criticamente os estudos selecionados na revisão da literatura que abordam, total ou parcialmente, o problema investigado. Mais do que descrever cada estudo, cabe compará-los segundo dimensões analíticas comuns (objetivo, método, contexto, dados e resultados), evidenciando o que já está estabelecido e o que permanece em aberto.

% Exemplo: substituir pelos estudos efetivamente selecionados no corpus da revisão.
O Quadro~\ref{qua:estudos_rsl} sintetiza as principais dimensões de comparação entre os estudos revisados e a proposta desta pesquisa, permitindo identificar semelhanças, diferenças e lacunas que motivam as escolhas metodológicas adotadas.

\begin{quadros}[H]
  \centering
  \caption{Síntese do corpus de estudos selecionados na revisão sistemática da literatura}
  \label{qua:estudos_rsl}
  \scriptsize
  \singlespacing
  \setlength{\tabcolsep}{2.5pt}
  \renewcommand{\arraystretch}{1.25}
  \begin{tabular}{|c|l|c|c|c|}
    \hline
    \textbf{ID}                                                                      &
    \textbf{Referência}                                                              &
    \textbf{Ano}                                                                     &
    \makebox[0.24\linewidth][c]{\textbf{Foco Metodológico}}                          &
    \makebox[0.38\linewidth][c]{\textbf{Principais Contribuições}}                        \\ \hline

    E01                                                                              &
    \citeonline{silva2024}                                                           &
    2024                                                                             &
    \parbox[t]{0.24\linewidth}{\raggedright
    Modelagem preditiva supervisionada com Random Forest e XGBoost.\strut}           &
    \parbox[t]{0.38\linewidth}{\raggedright
    Identificação de fatores críticos de risco espacial; validação cruzada k-fold.\strut} \\ \hline

    E02                                                                              &
    \citeonline{chen2024xai}                                                         &
    2024                                                                             &
    \parbox[t]{0.24\linewidth}{\raggedright
    Inteligência Artificial Explicável (XAI) com SHAP em dados espaciais.\strut}     &
    \parbox[t]{0.38\linewidth}{\raggedright
    Mapeamento de fidelidade local e estabilidade temporal dos preditores.\strut}         \\ \hline

    E03                                                                              &
    \citeonline{oliveira2025}                                                        &
    2025                                                                             &
    \parbox[t]{0.24\linewidth}{\raggedright
    Integração multissensorial de imagens orbitais Sentinel-2 e bases locais.\strut} &
    \parbox[t]{0.38\linewidth}{\raggedright
    Protocolos de calibração espectral e filtragem de ruído amostral.\strut}              \\ \hline

    E04                                                                              &
    \citeonline{santos2025}                                                          &
    2025                                                                             &
    \parbox[t]{0.24\linewidth}{\raggedright
    Calibração de limiares decisórios e sensibilidade paramétrica.\strut}            &
    \parbox[t]{0.38\linewidth}{\raggedright
    Otimização de revocação e controle de erro tipo II em classes raras.\strut}           \\ \hline

    E05                                                                              &
    \citeonline{martins2026}                                                         &
    2026                                                                             &
    \parbox[t]{0.24\linewidth}{\raggedright
    Métricas formais de interpretabilidade causal em redes profundas.\strut}         &
    \parbox[t]{0.38\linewidth}{\raggedright
    Proposição de índices quantitativos para governança e auditoria algorítmica.\strut}   \\ \hline
  \end{tabular}

  \fonte{Elaborado pelo autor (\ano).}
\end{quadros}

% ------------------------------------------------------------------------------
% EXEMPLO: tabela em modo paisagem (adaptar ou apagar)
% Use paisagem apenas quando a tabela não couber na largura da página em
% retrato, mesmo com fonte reduzida. O ambiente landscape (pacote lscape, já
% carregado pela classe) gira a página inteira e começa uma página nova; o
% texto que vem depois continua na página seguinte, em retrato.
% Regras da tabela (padrão IBGE, adotado pelo manual): dados numéricos,
% título acima, fonte abaixo, sem traços verticais e sem fechamento lateral;
% apenas os fios do cabeçalho e do rodapé (\toprule, \midrule, \bottomrule).
% Os valores abaixo são FICTÍCIOS, assim como as referências E01 a E05.
% ------------------------------------------------------------------------------
\begin{orientacao}[Tabela em modo paisagem]
  Quadro e tabela têm funções distintas. O quadro organiza informação textual e tem bordas em todos os lados, como o Quadro~\ref{qua:estudos_rsl}. A tabela apresenta dados numéricos, segue o padrão do IBGE e não tem traços verticais nem bordas laterais, como a Tabela~\ref{tab:desempenho_corpus}. A orientação paisagem é reservada a tabelas que não cabem na largura da página em retrato. A tabela deve ser citada no texto antes de aparecer, e os valores devem ser transcritos dos estudos originais, com a mesma métrica e o mesmo número de casas decimais.
\end{orientacao}

Além da comparação qualitativa, convém confrontar o desempenho que cada estudo reporta. A Tabela~\ref{tab:desempenho_corpus} reúne, para os mesmos estudos, o volume de dados, a proporção da classe positiva, o modelo de melhor resultado, o esquema de validação e as métricas de avaliação informadas pelos autores. A leitura conjunta dessas colunas mostra que valores elevados de AUC-ROC convivem com AUC-PR modesta nos estudos de baixa prevalência (E04 e E05), o que reforça a necessidade de métricas sensíveis ao desbalanceamento na avaliação de classes raras.

\begin{landscape}
  \begin{table}[H]
    \centering
    \caption{Desempenho preditivo reportado pelos estudos do corpus da revisão sistemática}
    \label{tab:desempenho_corpus}
    \footnotesize
    \singlespacing
    \setlength{\tabcolsep}{4pt}
    \renewcommand{\arraystretch}{1.2}
    \begin{tabular}{cllrrllccccl}
      \toprule
      \textbf{ID} &
      \textbf{Referência} &
      \makecell[l]{\textbf{Domínio}\\\textbf{dos dados}} &
      \makecell{\textbf{Observações}\\\textbf{(n)}} &
      \makecell{\textbf{Classe}\\\textbf{positiva (\%)}} &
      \makecell[l]{\textbf{Melhor}\\\textbf{modelo}} &
      \textbf{Validação} &
      \textbf{AUC-ROC} &
      \textbf{AUC-PR} &
      \textbf{F1} &
      \textbf{Revocação} &
      \textbf{XAI} \\
      \midrule
      E01 & \citeonline{silva2024}    & \makecell[l]{Riscos ambientais\\urbanos} & 18.450 & 12,4 & XGBoost       & \makecell[l]{$k$-fold\\($k=10$)} & 0,912 & 0,684 & 0,701 & 0,742 & Não              \\
      E02 & \citeonline{chen2024xai}  & \makecell[l]{Observação\\da Terra} & 52.300 &  8,7 & XGBoost       & \makecell[l]{Espacial\\(5 blocos)} & 0,887 & 0,611 & 0,645 & 0,690 & SHAP             \\
      E03 & \citeonline{oliveira2025} & \makecell[l]{Sentinel-2\\e cadastro} &  9.870 & 21,3 & \makecell[l]{Random\\Forest} & \makecell[l]{\textit{Holdout}\\70/30} & 0,856 & 0,703 & 0,672 & 0,705 & Não              \\
      E04 & \citeonline{santos2025}   & \makecell[l]{Anomalias\\territoriais} & 31.020 &  2,1 & LightGBM      & \makecell[l]{Estratificada\\($k=5$)} & 0,941 & 0,452 & 0,513 & 0,804 & Não              \\
      E05 & \citeonline{martins2026}  & \makecell[l]{Sistemas\\críticos} &  7.400 &  5,6 & \makecell[l]{Rede\\profunda} & \makecell[l]{\textit{Holdout}\\temporal} & 0,903 & 0,577 & 0,598 & 0,631 & \makecell[l]{Índices\\causais} \\
      \bottomrule
    \end{tabular}

    \fonte{Elaborada pelo autor (\ano), com base nos estudos citados.}
    {\footnotesize Nota: AUC-ROC: área sob a curva ROC; AUC-PR: área sob a curva de precisão e revocação. A AUC-PR depende da proporção da classe positiva e deve ser comparada com cautela entre estudos de prevalências distintas.\par}
  \end{table}
\end{landscape}

% ------------------------------------------------------------------------------
% EXEMPLO: tabela em paisagem GIRADA NO LEITOR DE PDF (adaptar ou apagar)
% Diferença em relação ao exemplo anterior: o ambiente "paisagem" (definido em
% config/preambulo.tex) também gira a página na tela, de modo que o leitor de
% PDF já a exibe na horizontal. Na impressão, os dois ambientes são iguais.
% Recursos ilustrados aqui:
%   \multicolumn{n}{c}{...}  título que agrupa n colunas
%   \cmidrule(lr){a-b}       fio parcial sob o grupo, das colunas a até b
%   ..                       sinal convencional do IBGE para "não se aplica"
% Os valores são FICTÍCIOS e coerentes com a Tabela de desempenho anterior.
% ------------------------------------------------------------------------------
\begin{orientacao}[Tabela girada para leitura na tela]
  Prefira o ambiente \texttt{paisagem} quando o trabalho for lido principalmente em tela, como na avaliação pela banca, pois a página já aparece na horizontal, sem que o leitor precise girá-la. Títulos de grupo (\texttt{\textbackslash multicolumn}) com fios parciais (\texttt{\textbackslash cmidrule}) evitam repetir o mesmo rótulo em várias colunas. Nas células sem valor, use os sinais convencionais do IBGE: ``..'' (não se aplica), ``...'' (dado não disponível) e ``--'' (valor igual a zero).
\end{orientacao}

O esquema de validação também condiciona a comparação. Os estudos que relataram resultados sob dois esquemas permitem estimar quanto do desempenho decorre da partição aleatória dos dados, que ignora a dependência espacial e temporal entre as observações. A Tabela~\ref{tab:validacao_corpus} contrapõe os dois cenários. Em todos os estudos, a validação que respeita essa dependência reduziu a AUC-PR, em média, em 8,1 pontos percentuais, e o F1 em 6,2 pontos. Os valores obtidos com partição aleatória tendem, portanto, a superestimar o desempenho esperado em operação.

\begin{paisagem}
  \begin{table}[H]
    \centering
    \caption{Desempenho reportado pelos estudos do corpus segundo o esquema de validação}
    \label{tab:validacao_corpus}
    \footnotesize
    \singlespacing
    \setlength{\tabcolsep}{5pt}
    \renewcommand{\arraystretch}{1.3}
    \begin{tabular}{cllcccccccc}
      \toprule
      & & &
      \multicolumn{3}{c}{\textbf{Validação aleatória}} &
      \multicolumn{3}{c}{\textbf{Validação espacial ou temporal}} &
      \multicolumn{2}{c}{\textbf{Variação (p.p.)}} \\
      \cmidrule(lr){4-6} \cmidrule(lr){7-9} \cmidrule(lr){10-11}
      \textbf{ID} & \textbf{Referência} & \textbf{Modelo} &
      \textbf{AUC-ROC} & \textbf{AUC-PR} & \textbf{F1} &
      \textbf{AUC-ROC} & \textbf{AUC-PR} & \textbf{F1} &
      \textbf{AUC-PR} & \textbf{F1} \\
      \midrule
      E01 & \citeonline{silva2024}    & XGBoost       & 0,912 & 0,684 & 0,701 & 0,861 & 0,592 & 0,633 & $-9{,}2$ & $-6{,}8$ \\
      E02 & \citeonline{chen2024xai}  & XGBoost       & 0,934 & 0,702 & 0,711 & 0,887 & 0,611 & 0,645 & $-9{,}1$ & $-6{,}6$ \\
      E03 & \citeonline{oliveira2025} & Random Forest & 0,856 & 0,703 & 0,672 & 0,812 & 0,641 & 0,618 & $-6{,}2$ & $-5{,}4$ \\
      E04 & \citeonline{santos2025}   & LightGBM      & 0,941 & 0,452 & 0,513 & 0,902 & 0,381 & 0,462 & $-7{,}1$ & $-5{,}1$ \\
      E05 & \citeonline{martins2026}  & Rede profunda & 0,948 & 0,668 & 0,671 & 0,903 & 0,577 & 0,598 & $-9{,}1$ & $-7{,}3$ \\
      \midrule
      \multicolumn{3}{l}{\textbf{Média}} & .. & .. & .. & .. & .. & .. & $\mathbf{-8{,}1}$ & $\mathbf{-6{,}2}$ \\
      \bottomrule
    \end{tabular}

    \fonte{Elaborada pelo autor (\ano), com base nos estudos citados.}
    {\footnotesize Nota: Variação em pontos percentuais (p.p.) entre a validação espacial ou temporal e a aleatória. E05 adota validação temporal; os demais, espacial. Sinal convencional: .. não se aplica.\par}
  \end{table}
\end{paisagem}

O fechamento desta seção deve retomar, de forma argumentativa, a lacuna anunciada na Seção~\ref{sec:problema}. É nesse ponto que a comparação entre os estudos deixa de ser descritiva e passa a justificar, com base na própria literatura, a pertinência e a originalidade da investigação.
